\documentclass[11pt,a4paper]{article}
\usepackage[margin=23mm]{geometry}
\usepackage{amsmath,amssymb,array,booktabs}
\input{glyphtounicode}
\pdfgentounicode=1
\usepackage[hidelinks]{hyperref}
\hypersetup{hypertexnames=false}
\setlength{\parindent}{0pt}
\setlength{\parskip}{5pt}
\newcommand{\pos}[1]{[#1]_+}
\renewcommand{\labelitemi}{\(\bullet\)}
\begin{document}
\begin{center}
{\Large R-E selected: boundaries, recovery overlap and the next choice}\\[4pt]
15 October 2026 --- analytical review; no traffic time stepping
\end{center}

\textbf{Status.} The user selects R-E, closing SEP-R1 at source-attribution
scope. Ordinary A departures do not add independent recovery profiles.
S-E-P remains selected. Exact work, marks, overlap and global event proofs
remain open. This continuation stops at \textbf{RE-O1}, Type 2 G3: additive
burdens O-S versus recent-mass renewal O-R. Neither overlap law is selected.

\section*{1. Retained interface and the trial isolated profile}

Retain H1's entrance histories $F_j$, advice $v\in[v_-,v_+]$, finite
terminal length $b$, initially empty histories, and
\begin{equation}
 P_j(t)=F_j(t-b/v),\quad q_j=P_j-Y_j,\quad N_j=F_j-Y_j,
 \quad \lambda_j=\dot P_j,\quad g_j=\dot Y_j. \tag{1}
\end{equation}
The B lane remains independent and unqueued; $B(t)$ below is instead the
shared A/E budget. Write $K=wk_j$ and use
\begin{equation}
 c(u)=\frac{Ku}{w+u},\quad C(v)=c(v),\quad
 U(f)=\frac{wf}{K-f},\quad 0<R\leq C(v_-),\quad u_R=U(R). \tag{2}
\end{equation}
The trial node has ordinary cost one, $c_0=1+\delta$, and
\begin{equation}
 c_E=c_0+\beta\frac{g_A}{C(v)}\frac{\pos{u_A-u_E}}v,
 \quad g_A+c_Eg_E\leq B,\quad
 \eta=c_0-\beta(1+v_+/w)>0. \tag{3}
\end{equation}
Here $\delta,\beta\geq0$. The sufficient $\eta$ condition and representative
trace marks $u_j=v$ when free and $u_j=U(g_j)$ when queued remain trial
assumptions. The earlier note proves selected-regime roots, not a complete
hybrid law. S-E-P protects a free A offer that fits $B$; queued A uses the
equal-work/unused-work redistribution template, even when E is free.

For $a_r>0$, the proposed isolated recovery profile is
\begin{align}
 V(\theta;u)&=\sqrt{(w+u)^2+2wa_r\theta}-w,\nonumber\\
 D(\theta;u)&=\pos{R-c(V(\theta;u))},\qquad \theta\geq0. \tag{4}
\end{align}
For $u<u_R$, its duration and integrated loss are
\begin{equation}
 \Theta(u)=\frac{u_R-u}{a_r}+\frac{u_R^2-u^2}{2a_rw},\quad
 W(u)=\frac{k_c(u_R)(u_R-u)^2}{2a_r},\quad
 k_c(z)=\frac{K}{w+z}. \tag{5}
\end{equation}
For $u\geq u_R$ both are zero. Put $\Theta_{\max}=\Theta(0)$,
$W_0=W(0)$. Then $0\leq D\leq R$, $D$ decreases with age,
$|D_\theta|\leq a_rk_j$, and $|D_u|\leq k_j$.
The age includes backward propagation, not a numerical time step.

\newpage
\section*{2. What the additive R-E law now proves}

O-S is the existing trial overlap family, with physical reference mass
$m=m_{\rm ref}=1$ vehicle-equivalent:
\begin{equation}
 p_S(t)=\frac1m\int_0^tD(t-s;u_E(s))g_E(s)ds,\qquad
 B_S(t)=\pos{R-p_S(t)}. \tag{6}
\end{equation}
This is one account. It does not create a vehicle queue or another
constraint for each past passage. The initial account is neutral.

\textbf{Strictly positive finite-time service, conditional on an admissible
causal solution.} Since $c_E\geq c_0$, work feasibility gives
$g_E\leq B_S/c_0$. Differentiating the nonincreasing profiles gives,
almost everywhere,
\begin{equation}
 \dot p_S=\frac{g_E(t)D(0;u_E(t))}m+
  \frac1m\int_0^tD_\theta(t-s;u_E(s))g_E(s)ds
 \leq\frac Rm g_E\leq\frac{R}{mc_0}B_S. \tag{7}
\end{equation}
Before a hypothetical first zero of $B_S$, $B_S=R-p_S$ and
\begin{equation}
 \dot B_S\geq-\frac{R}{mc_0}B_S,\qquad
 B_S(t)\geq R\exp[-Rt/(mc_0)]>0. \tag{8}
\end{equation}
Continuity rules out a finite first zero. This derives a bound; it does not
insert a positive floor. An arbitrary externally prescribed admission
history can exhaust (6); it need not be admissible in the coupled node.
Neither (8) nor the earlier fixed-mode contraction proves global queue-event
existence, uniqueness or empirical fidelity.

\textbf{A reachable first loss, not just a possible source.} At the retained
H1 fast-front example, translated so first eligibility is $t=0$,
\begin{equation}
 R=39/20,\quad \lambda_A=1047/880,\quad \lambda_E=125/176,
 \quad \delta=\beta=1/10,\quad C=3. \tag{9}
\end{equation}
The earlier strict-onset calculation yields free A, queued E and
\begin{equation}
 g_* =\frac{16716-183\sqrt{5234}}{5189},\quad
 \frac{2230}{3343}<g_*<\frac{669}{968}<\lambda_E<R. \tag{10}
\end{equation}
Thus $u_E=U(g_*)<u_R$ and, with empty history,
\begin{equation}
 \dot p_S(0+)=\frac{g_*(R-g_*)}m>0,\qquad
 \dot B_S(0+)=-\frac{g_*(R-g_*)}m<0. \tag{11}
\end{equation}
The local fixed-mode causal construction and strict onset margin justify
this initial right-hand effect. These numbers are inherited diagnostic
values, not selected task calibration or a positive-control parameter set.
No complete A-delay witness, preferred advice or interior optimum follows.
The O-R candidate below agrees until the first reference mass has passed.

\newpage
\section*{3. Strict boundary compatibility and a benign full-task reduction}

Write $a=\lambda_A$, $e=\lambda_E$. With queued E, its unique trial
service solves $H_E(g_E)=B-g_A$ and satisfies
$g_E\leq(B-g_A)/c_0$. With both queues active, $g_A=B/2$.
With free E and queued A, E's entire offer fits if $c_0e\leq B/2$;
then $g_A=B-c_0e$. These are current-state rules, not allocation latches.

\begin{center}
\small
\begin{tabular}{p{0.22\linewidth}p{0.30\linewidth}p{0.38\linewidth}}
\toprule
Strict event & Condition before switch & Compatible right-hand behavior\\
\midrule
E onset; A free & $a\leq B$, $e>(B-a)/c_0$ &
Queued E has $g_E\leq(B-a)/c_0<e$, so its newborn queue grows.\\[4pt]
A onset & $a>B$ & Queued allocation has $g_A\leq B<a$;
if free E also fails its share, its queued root is below that free share.\\[4pt]
E onset; A queued & $e>B/(2c_0)$ &
Queued E has $g_E\leq B/(2c_0)<e$.\\[4pt]
E clears; A free & $e<g_E^Q\leq(B-a)/c_0$ &
Both free offers fit $a+c_0e<B$.\\[4pt]
E clears; A queued & $e<g_E^Q\leq B/(2c_0)$ &
Free E gets $e$; A gets $B-c_0e\geq B/2$. Cooperation persists.\\[4pt]
A clears; E queued & $a<g_A^Q=B/2$ &
Free A's offer fits $B$; E is recomputed against its residual.\\[4pt]
A clears; E free & $a<g_A^Q=B-c_0e$ &
Both free offers fit $a+c_0e<B$.\\
\bottomrule
\end{tabular}
\end{center}
For simultaneous strict clearance of both queues, $a<B/2$ and
$e<g_E^Q\leq B/(2c_0)$ imply free admission is feasible. These statements
check switch directions. They do not establish every grazing or simultaneous
birth case, an absence of Zeno accumulation, or full hybrid uniqueness.
The earlier no-spontaneous-birth guard remains required. A step size,
queue epsilon, sliding convention or extra hysteresis has not been added.

\textbf{Exact benign reduction.} For one H1 advice, translate its eligible
pulse to $0<t<\ell/v$. On that interval
\begin{equation}
 a(t)+e(t)=s:=v(e_0+a_0)<R,\quad e(t)\ \text{nonincreasing},\quad
 a(t)+c_0e(t)=s+\delta e(t). \tag{12}
\end{equation}
If $s+\delta e(0+)\leq R$, the free offers fit from the outset and only
become easier to serve. Free E has $u_E=v\geq u_R$, so $D=0$ and no new
account loss arises. Causal free-mode retention therefore gives
\begin{equation}
 q_A=q_E=0,\quad p=0,\quad B=R,\quad
 J(v)=M(L-\ell/2)/v. \tag{13}
\end{equation}
If this holds for every admitted advice, all observers have the fastest
optimum. This is a declared exact negative case for either overlap candidate,
not a reason to force queueing or choose parameters for a positive result.

\newpage
\section*{4. A stronger simple rival and the overlap reference}

\textbf{Stationary-equivalent work.} If $g_E$ and its mark $u_E$ have
been constant over a full profile window, O-S gives
\begin{equation}
 p_S=\frac{g_EW(u_E)}m,\qquad
 g_A+\left(c_E+\frac{W(u_E)}m\right)g_E=R
 \quad\text{at saturated service}. \tag{14}
\end{equation}
A state-dependent instantaneous work relation therefore matches this
stationary consequence. This is our deduction, not a new native invariant
junction. Its transient selection and feasibility need their own contract.
It is a serious simpler rival: a history benefit must concern consequential
transients, not merely lower average discharge. Proper aggregate/history,
scalar-statistic, native lane-aware and equal-information search rivals stay.

\textbf{Point-reference screen.} The homogeneous point case in
\href{https://repository.tudelft.nl/file/File_c424c1b3-7374-436b-a6a7-abdf69063aa6}{Leclercq et al. (2016), Section 2, equation (1)}
evaluates cumulative passing over the interval to the next insertion.
For a fixed mark $u<u_R$ and fixed headway $h_0$, capping our isolated
service at $R$ gives the following own interval-mean deduction:
\begin{equation}
 \overline B_{\rm point}=R-\frac{W_{h_0}(u)}{h_0},\qquad
 W_h(u):=\int_0^hD(\theta;u)d\theta. \tag{15}
\end{equation}
For $z_h=\min\{u_R,V(h;u)\}$,
\begin{equation}
 W_h(u)=\frac{k_c(u_R)}{2a_r}
       \{(u_R-u)^2-(u_R-z_h)^2\}. \tag{16}
\end{equation}
Full-profile addition under the corresponding constant mean admission
$g_E=m/h_0$ instead gives $\overline p_S=W(u)/h_0$. Hence
\begin{equation}
 \overline p_S-\frac{W_{h_0}(u)}{h_0}
   =\frac1{h_0}\int_{h_0}^{\Theta(u)}D(\theta;u)d\theta>0
 \quad(0<h_0<\Theta(u)). \tag{17}
\end{equation}
When $h_0\geq\Theta(u)$ the difference is zero. This compares an additive
mean closure with a capped homogeneous interval reference. It is not an
exact native transient law, a primary H1 trajectory, or a proof that the
closed-loop O-S model exhausts its budget. Equation (8) prevents the latter
at finite time in an admissible initially neutral solution.

An exact constitutive diagnostic uses $w=2,k_j=3,R=2,u=1/2,a_r=1/2$,
$m=h_0=1$. It gives $u_R=1$, $\Theta=11/8$, and
\begin{equation}
 W=1/2,\quad W_1=6\sqrt{33}-34,\quad
 W-W_1=69/2-6\sqrt{33}>0. \tag{18}
\end{equation}
The mark and mean admissions are prescribed for this screen, not asserted
to solve the trial node. No traffic is time stepped. Variable marks and
interacting restrictions still require a physical overlap closure.

\newpage
\section*{5. A concrete recent-mass renewal alternative}

For O-R, let $T_E(\xi)$ be a generalized inverse of the continuous
admitted history $Y_E$. Endpoint choices at plateaus have zero mass measure.
Keep the same profile and define
\begin{equation}
 p_R(t)=\frac1m\int_{\pos{Y_E(t)-m}}^{Y_E(t)}
     D\bigl(t-T_E(\xi);u_E(T_E(\xi))\bigr)d\xi,
 \qquad B_R=\pos{R-p_R}. \tag{19}
\end{equation}
Equivalently integrate (6) only over past mass with
$0\leq Y_E(t)-Y_E(s)<m$. This uses predicted continuous mass, not integer
vehicles, packets, a second traffic queue or initial hidden information.
Only still-active profiles within the most recent reference mass contribute.
At first admission and until $Y_E=m$, O-R and O-S give the same functional.
Their later closed-loop trajectories need not be ordered.

The physical interpretation is renewal of an average affected stream.
An admitted portion contributes while replacing an equal amount of oldest
profile weight after the window is full. The same $m=1$ reference unit now
also sets a physical renewal mass. Its role must be defended, not changed
as a harmless numerical resolution or tuned for a desired control result.

\textbf{Bounds that follow immediately.} For any fixed admissible history,
\begin{equation}
 0\leq p_R(t)\leq p_S[Y_E,u_E](t),\qquad
 p_R(t)\leq R\min\{Y_E(t)/m,1\}. \tag{20}
\end{equation}
The comparison uses the same supplied history; it does not compare two
endogenously different solutions. After last E admission, the retained
window is fixed and its profiles decrease, expiring within
$\Theta_{\max}$. A later fast E portion has zero new profile but can retire
older slow weight. That is part of this renewal model and a material risk:
it can release service too quickly when an older leader still controls it.

For fixed $g_E>0$, mark $u$, and enough previous admissions,
\begin{equation}
 h_0=m/g_E,\qquad
 p_R=\frac{g_E}{m}\int_0^{h_0}D(\theta;u)d\theta
      =\frac{W_{h_0}(u)}{h_0}. \tag{21}
\end{equation}
Thus O-R reproduces (15)'s mean loss in that restricted reference. Equation
(19) is our transient extension, not the source's native algorithm. Its
window cutoff needs its own causal/event qualification. The O-S inverse
and fixed-mode history proof do not automatically establish O-R uniqueness.

Both candidates preserve one shared output and E-only new disturbance.
O-S preserves every older physical aging tail but assumes independent
addition. O-R limits repeated accounting but assumes recent-mass renewal
adequately represents old restrictions. No new spatial field, longitudinal
zone, approach-speed ODE or A-source penalty is introduced by either option.

\newpage
\section*{6. Joined accounting, proof limits and stopping gate}

For either overlap law, let $\overline M_E=P_E(L/v)$ and assume an
admissible causal solution. Nonnegative kernels and (20) give
\begin{equation}
 \int_0^\infty p(t)dt\leq\frac{\overline M_EW_0}{m},\qquad
 R-B=\min\{R,p\}\leq p. \tag{22}
\end{equation}
Using the prior work-conserving trial node, $c_{\max}=c_0+\beta$, and
last eligibility $t_*=L/v$, the existing clearance argument gives
\begin{equation}
 T_{\rm clear}\leq t_*+
 \frac{c_{\max}(q_A(t_*)+q_E(t_*))+\overline M_EW_0/m}{R}. \tag{23}
\end{equation}
It uses an upper bound on available E mass before clearance, not an assumed
successful final exit. This is a conditional bound, not proof of a full
hybrid solution. No positive service floor is added.

For $M<k_jb$ and $\dot F_j\leq C(v)$, the earlier finite-terminal vacancy
proof is unchanged for nonnegative departures. Physical population and
the complete local objective remain
\begin{equation}
 s(t)\rho_0+\sum_jN_j+\sum_jY_j=M,\qquad
 J(v)=\frac{M(L-\ell/2)}v+\int_0^\infty(q_A+q_E)dt. \tag{24}
\end{equation}
Early-merged mass belongs to A, not a second forced E merge. A retired
recovery profile is not deleted traffic. Post-line recovery memory is not
post-line vehicle-time. The independent B lane has no waiting in this domain.

\textbf{State and evidence.} Current allocation uses current traffic and
budget. Future recovery generally needs admission/speed history, not just
the scalar reading. Under initial symmetry and known empty memory, total
mass plus initial ending mass still reconstructs H1; unique reconstruction
of the joined dynamics remains conditional on a qualified whole law.
No full-lane-state necessity, new action split, gradient correctness or
independent gradient usefulness is proved. The constant-A protection result
remains; changing H1 arrivals need a consequential complete A-delay/MP6
witness. A positive local loss derivative in (11) does not supply that.

\textbf{Source and scope.} Carey's Section 5 own-lane/work ingredients and
Jin's boundary consistency warning were rechecked at prior reading scope;
neither supplies this trial node's full event proof. S13, the closed urban
task, the closed ramp line and the model--SUMO S-P--SUMO S-L ladder retain
their status. No reference simulator, fitting or numerical gradients ran.

\textbf{Gate RE-O1.} Select independent additive fading burdens O-S, or
recent-mass renewal O-R. R-E source attribution stays closed and S-E-P stays
selected. Provisionally qualify O-R first for its explicit overlap treatment
and restricted point-reference correspondence, retaining O-S as a serious
phenomenological rival. Both are concrete formulas with different physical
assumptions; neither is selected. Stop before adopting the new overlap law.
The approved joined pass and Step 3 remain in progress; Step 4 is unopened.
\end{document}
